\documentclass[11pt]{scrartcl}
\usepackage[utf8]{inputenc}
\usepackage[sexy]{evanchen} 
\usepackage{graphicx}
\graphicspath{{8 Outsource K Math/Gambar/}}
\usepackage{multicol}


\title{Trigonometri}
\author{Azzam and Akeyla}
\date{August 2022}

\begin{document}

\maketitle

\section{Persamaan}
Tentunya, banyak properti lain yang belum dicover modul ini. Namun, sebagian besar properti yang berguna akan ditulis di sini.
\subsection{Well-Known}
\begin{enumerate}
    \item $\sin^2{x} + \cos^2{x} = 1$
    \item $\sin{(x \pm y)} = \sin{x}\cos{y} \pm \sin{y}\cos{x}$
    \item $\cos{(x \pm y)} = \cos{x}\cos{y} \mp \sin{x}\sin{y}$
    \item $\tan{(x + y)} = \dfrac{\tan{x} \pm \tan{y}}{1 \mp \tan{x}\tan{y}}$
\end{enumerate}
\subsection{Double Angle}
\begin{enumerate}
    \item $\sin{2x} = 2\sin{x}\cos{x}$
    \item $\cos{2x} = \cos^2{x} - \sin^2{x} = 2\cos^2{x} - 1 = 1 - 2\sin^2{x}$
    \item $\tan{2x} = \dfrac{2\tan{x}}{1 - \tan^2{x}}$
\end{enumerate}
\subsection{Triple Angle}
\begin{enumerate}
    \item $\sin{3x} = 3\sin{x} - 4\sin^3{x}$
    \item $\cos{3x} = 4\cos^3{x} - 3\cos{x}$
    \item $\tan{3x} = \dfrac{3\tan{x} - \tan^3{x}}{1 - 3\tan^2{x}}$
\end{enumerate}
\subsection{Sum to Product}
\begin{enumerate}
    \item $\sin{x} + \sin{y} = 2\sin{\left(\dfrac{x + y}{2}\right)}\cos{\left(\dfrac{x - y}{2}\right)}$
    \item $\cos{x} + \cos{y} = 2\cos{\left(\dfrac{x + y}{2}\right)}\cos{\left(\dfrac{x - y}{2}\right)}$
\end{enumerate}
\subsection{Product to Sum}
\begin{enumerate}
    \item $\sin{x}\sin{y} = \dfrac{1}{2}(\cos{(x - y)} - \cos{(x + y)})$
    \item $\cos{x}\cos{y} = \dfrac{1}{2}(\cos{(x + y)} + \cos{(x - y)})$
    \item $\sin{x}\cos{y} = \dfrac{1}{2}(\sin{(x + y)} - \sin{(y - x)})$
\end{enumerate}
\subsection{Half tangent}
Misal $t = \tan{\frac{x}{2}}$. Maka,
\begin{enumerate}
    \item $\sin{x} = \dfrac{2t}{1 + t^2}$
    \item $\cos{x} = \dfrac{1 - t^2}{1 + t^2}$
    \item $\tan{x} = \dfrac{2t}{1 - t^2}$
\end{enumerate}
\subsection{Trigonometri Sudut pada Segitiga}
\label{subsec:triangle_trigs}
Misal diberikan sebuah segitiga $ABC$ dengan panjang sisi $BC = a, CA = b,$ dan $AC = b$ serta besar sudut $\angle{BAC} = \alpha, \angle{CBA} = \beta$, dan $\angle{ABC} = \beta$. Beberapa properti:
\begin{enumerate}
    \item $\sin{\alpha} + \sin{\beta} + \sin{\gamma} = 4\cos{\frac{\alpha}{2}}\cos{\frac{\beta}{2}}\cos{\frac{\gamma}{2}}$
    \item $\cos{\alpha} + \cos{\beta} + \cos{\gamma} = 1 + 4\sin{\frac{\alpha}{2}}\sin{\frac{\beta}{2}}\sin{\frac{\gamma}{2}}$
    \item $\tan{\alpha}\tan{\beta}\tan{\gamma} = \tan{\alpha} + \tan{\beta} + \tan{\gamma}$
    \item $\tan{\frac{\alpha}{2}}\tan{\frac{\beta}{2}} + \tan{\frac{\beta}{2}}\tan{\frac{\gamma}{2}} + \tan{\frac{\gamma}{2}}\tan{\frac{\alpha}{2}} = 1$
    \item $\sin{2\alpha} + \sin{2\beta} + \sin{2\gamma} = 4\sin{\alpha}\sin{\beta}\sin{\gamma}$
    \item $\cos{2\alpha} + \cos{2\beta} + \cos{2\gamma} = -1 - 4\cos{\alpha}\cos{\beta}\cos{\gamma}$
    \item $sin^2{\alpha} + \sin^2{\beta} + \sin^2{\gamma} = 2 + 2\cos{\alpha}\cos{\beta}\cos{\gamma}$
    \item $\sin^2{\frac{\alpha}{2}} + \sin^2{\frac{\beta}{2}} + \sin^2{\frac{\gamma}{2}} + 2\sin{\frac{\alpha}{2}}\sin{\frac{\beta}{2}}\sin{\frac{\gamma}{2}} = 1$
    \item $\cos^2{\alpha} + \cos^2{\beta} + \cos^2{\gamma} + 2\cos{\alpha}\cos{\beta}\cos{\gamma} = 1$
\end{enumerate}

\subsection{Jari-jari lingkaran luar dan dalam}
Notasikan $R$ dan $r$ secara berturut-turut sebagai panjang jari-jari lingkaran luar dan lingkaran dalam
\begin{enumerate}
    \item (Dalil sinus) $\dfrac{a}{\sin{\alpha}} = \dfrac{b}{\sin{\beta}} = \dfrac{c}{\sin{\gamma}} = 2R$
    \item $2R\sin{\alpha}\sin{\beta}\sin{\gamma} = r(\sin{\alpha} + \sin{\beta} + \sin{\gamma})$
    \item $r = 4R\sin{\frac{\alpha}{2}}\sin{\frac{\beta}{2}}\sin{\frac{\gamma}{2}}$
\end{enumerate}

\subsection{Soal-soal}
\begin{problem}
    Buktikan bahwa \[\dfrac{\tan{3x}}{\tan{x}} = \tan{\left(\dfrac{\pi}{3} - x\right)}\tan{\left(\dfrac{\pi}{3} + x\right)}\] untuk setiap $x \neq \frac{k\pi}{6}$ dengan $k \in \mathbb{Z}$.
    % 103 29
\end{problem}

\begin{problem}
    (Turki MO 1998) Barisan $\{a_n\}$ didefinisikan sebagai $a_1 = t$ dan $a_{n+ 1} = 4a_n(1 - a_n)$ untuk setiap $n \geq 1$. Tentukan banyak bilangan real $t$ berbeda yang bisa dipilih agar $a_{1998} = 0$
    % 103 20
\end{problem}

\begin{problem}
    Tentukan semua nilai $x$ yang memenuhi persamaan \[\sin{x} + \sin{2x} + \sin{3x} = 0\]
    % Krechmar Algebra 4 53
\end{problem}

\begin{problem}
    Tentukan semua nilai $x$ yang memenuhi persamaan \[\sin{5x} = 16 \sin^5{x}\]
    % Krechmar Algebra 4 54
\end{problem}

\begin{problem}
    Tentukan semua nilai $x$ yang memenuhi persamaan \[(1 - \tan{x})(1 + \sin{2x}) = 1 + \tan{x}\]
    % Krechmar Algebra 4 59
\end{problem}

\begin{problem}
    Buktikan bahwa jika \[\tan{x} + \tan{2x} + \tan{3x} + \tan{4x} = 0\] maka $5x = k\pi$ atau $8 \cos{2x} = 1\pm \sqrt{17}$
    % Krechmar Algebra 4 60
\end{problem}

\section{Ketaksamaan Trigonometri}
Notasikan $\alpha, \beta, \gamma$ sebagai tiga sudut dari segitiga $ABC$.
\subsection{Jensen-able atau AM-GM-able}
\begin{multicols}{2}
\begin{enumerate}
    \item $\sin{\alpha} + \sin{\beta} + \sin{\gamma} \leq \frac{3\sqrt{3}}{2}$
    \item $\sin{\alpha}\sin{\beta}\sin{\gamma} \leq \dfrac{3\sqrt{3}}{8}$
    \item $\sin{\frac{\alpha}{2}} + \sin{\frac{\beta}{2}} + \sin{\frac{\gamma}{2}} \leq \frac{3}{2}$
    \item $\sin{\frac{\alpha}{2}}\sin{\frac{\beta}{2}}\sin{\frac{\gamma}{2}} \leq \frac{1}{8}$
    \item $\cos{\frac{\alpha}{2}} + \cos{\frac{\beta}{2}} + \cos{\frac{\gamma}{2}} \leq \frac{3\sqrt{3}}{2}$
    \item $\cos{\frac{\alpha}{2}}\cos{\frac{\beta}{2}}\cos{\frac{\gamma}{2}} \leq \frac{3\sqrt{3}}{8}$
    \item $\tan{\frac{\alpha}{2}} + \tan{\frac{\beta}{2}} + \tan{\frac{\gamma}{2}} \geq \sqrt{3}$
    \item $\tan{\alpha} + \tan{\beta} + \tan{\gamma} \geq 3\sqrt{3}$
\end{enumerate}
\end{multicols}
\subsection{Manfaatkan subsection \ref{subsec:triangle_trigs}}
\begin{multicols}{2}
\begin{enumerate}   
    \item $\cos{\alpha} + \cos{\beta} + \cos{\gamma} \leq \frac{3}{2}$
    \item $\cos{\alpha}\cos{\beta}\cos{\gamma} \leq \frac{1}{8}$
    \item $\sin^2{\alpha} + \sin^2{\beta} + \sin^2{\gamma} \leq \frac{9}{4}$
    \item $\cos^2{\alpha} + \cos^2{\beta} + \cos^2{\gamma} \geq \frac{3}{4}$
    \item $\sin^2{\frac{\alpha}{2}} + \sin^2{\frac{\beta}{2}} + \sin^2{\frac{\gamma}{2}} \geq \frac{3}{4}$
    \item $\cos^2{\frac{\alpha}{2}} + \cos^2{\frac{\beta}{2}} + \cos^2{\frac{\gamma}{2}} \leq \frac{9}{4}$
    \item $\sin{2\alpha} + \sin{2\beta} + \sin{2\gamma} \leq \frac{3\sqrt{3}}{2}$
    \item $\cos{2\alpha} + \cos{2\beta} + \cos{2\gamma} \geq -\frac{3}{2}$
\end{enumerate}
\end{multicols}

\subsection{Soal-soal}
\begin{problem}
    Diberikan $0 \leq \alpha, \beta, \gamma \leq 90^{\circ}$ yang memenuhi $\sin{\alpha} + \sin{\beta} + \sin{\gamma} = 1$. Buktikan ketaksamaan \[\tan^2{\alpha} + \tan^2{\beta} + \tan^2{\gamma} \geq \dfrac{3}{8}\]
    % cvetkovski 263
\end{problem}

\begin{problem}
    Misal $\alpha, \beta, \gamma$ adalah sudut dari sebuah segitiga. Buktikan ketaksamaan \[\sin{\alpha} + \sin{\beta} + \sin{\gamma} \geq \sin{2\alpha} + \sin{2\beta} + \sin{2\gamma}\]
    % cvetkovski 260
\end{problem}

\begin{problem}
    Tentukan nilai minimal dari \[\dfrac{\sec^4{\alpha}}{\tan^2{\beta}} + \dfrac{\sec^4{\beta}}{\tan^2{\alpha}}\] untuk setiap bilangan real $\alpha, \beta \neq \frac{k\pi}{2}$ dengan $k \in \mathbb{Z}$
    % 103 39
\end{problem}

\begin{problem}
    Buktikan ketaksamaan \[\dfrac{\sin^3{a}}{\sin{b}} + \dfrac{\cos^3{a}}{\cos{b}} \geq \sec{(a - b)}\] untuk setiap $0 < a, b < \frac{\pi}{2}$
    % 103 43
\end{problem}

\begin{problem}
    Tentukan semua bilangan real $x$ yang memenuhi ketaksamaan $0 < x < \pi$ dan $\dfrac{8}{3 \sin x - \sin 3x} + 3\sin^2 x \leq 5$.
    % Vietnam MO 1969 P2
\end{problem}

\begin{problem}
    Misal $\tan{\theta} = n \tan{\varphi}$ dengan $n$ merupakan bilangan real positif. Buktikan bahwa \[\tan^2{(\theta - \varphi)} \leq \dfrac{(n - 1)^2}{4n}\]
    % Krechmar 8 10
\end{problem}

\begin{problem}
    Buktikan bahwa jika \[\dfrac{1}{\cos{\alpha}\cos{\beta}}+ \tan{\alpha}\tan{\beta} = \tan{\gamma}\] maka $\cos{2\gamma} \leq 0$.
    % Krechmar 8 11
\end{problem}

\section{Substitusi Trigonometri pada Ketaksamaan}
\subsection{Bentuk substitusi yang umum ditemui}
\begin{enumerate}
    \item Jika $a, b, c$ adalah bilangan real positif yang memenuhi $ab + bc + ca = 1$, maka terdapat sebuah segitiga dengan besar sudut $\alpha, \beta, \gamma$ ($\alpha + \beta + \gamma = \pi$) yang memenuhi \[a = \tan{\dfrac{\alpha}{2}}, b = \tan{\dfrac{\beta}{2}}, c = \tan{\dfrac{\gamma}{2}}\]
    \item Jika $a, b, c$ adalah bilangan real positif yang memenuhi $ab + bc + ca = 1$, maka terdapat sebuah segitiga \textbf{lancip} dengan besar sudut $\alpha, \beta, \gamma$ ($\alpha + \beta + \gamma = \pi$, $0 < \alpha, \beta, \gamma < \frac{\pi}{2}$) yang memenuhi \[a = \cot{\alpha}, b = \cot{\beta}, c = \cot{\gamma}\]
    \item Jika $a, b, c$ adalah bilangan real positif yang memenuhi $a + b + c = abc$, maka terdapat sebuah segitiga dengan besar sudut $\alpha, \beta, \gamma$ ($\alpha + \beta + \gamma = \pi$) yang memenuhi \[a = \cot{\dfrac{\alpha}{2}}, b = \cot{\dfrac{\beta}{2}}, c = \cot{\dfrac{\gamma}{2}}\]
    \item Jika $a, b, c$ adalah bilangan real positif yang memenuhi $a + b + c = abc$, maka terdapat sebuah segitiga \textbf{lancip} dengan besar sudut $\alpha, \beta, \gamma$ ($\alpha + \beta + \gamma = \pi$, $0 < \alpha, \beta, \gamma < \frac{\pi}{2}$) yang memenuhi \[a = \tan{\alpha}, b = \tan{\beta}, c = \tan{\gamma}\]
    \item Jika $a, b, c$ adalah bilangan real positif yang memenuhi $a^2 + b^2 + c^2 + 2abc = 1$, maka terdapat sebuah segitiga \textbf{lancip} dengan besar sudut $\alpha, \beta, \gamma$ ($\alpha + \beta + \gamma = \pi$, $0 < \alpha, \beta, \gamma < \frac{\pi}{2}$) yang memenuhi \[a = \sin{\dfrac{\alpha}{2}}, b = \sin{\dfrac{\beta}{2}}, c = \sin{\dfrac{\gamma}{2}}\]
    \item Jika $a, b, c$ adalah bilangan real positif yang memenuhi $a^2 + b^2 + c^2 + 2abc = 1$, maka terdapat sebuah segitiga \textbf{lancip} dengan besar sudut $\alpha, \beta, \gamma$ ($\alpha + \beta + \gamma = \pi$, $0 < \alpha, \beta, \gamma < \frac{\pi}{2}$) yang memenuhi \[a = \cos{\alpha}, b = \cos{\beta}, c = \cos{\gamma}\]
\end{enumerate}
\subsection{Soal-soal}
\begin{problem}
    Diberikan bilangan real positif $a, b, c$ yang memenuhi $a + b + c = abc$. Buktikan ketaksamaan \[\dfrac{1}{\sqrt{1 + a^2}} + \dfrac{1}{\sqrt{1 + b^2}} + \dfrac{1}{\sqrt{1 + c^2}} \leq \dfrac{3}{2}\]
    % cvetkivski 210
\end{problem}

\begin{problem}
    $a, b, c$ adalah bilangan real positif yang memenuhi $a + b + c = 1$. Buktikan ketaksamaan \[\sqrt{\dfrac{1}{a} - 1}\sqrt{\dfrac{1}{b} - 1} + \sqrt{\dfrac{1}{b} - 1}\sqrt{\dfrac{1}{c} - 1} + \sqrt{\dfrac{1}{c} - 1}\sqrt{\dfrac{1}{a} - 1} \geq 6\]
    % cvetkovski 213
\end{problem}

\begin{problem}
    $a, b, c$ adalah bilangan real nonnegatif yang memenuhi $ab + bc + ca = 1$. Buktikan ketaksamaan \[\dfrac{a}{1 + a^2} + \dfrac{b}{1 + b^2} + \dfrac{c}{1 + c^2} \leq \dfrac{3\sqrt{3}}{4}\]
    % cvetkovski 215
\end{problem}


\begin{problem}
    (Romania MO 2002) Jika $a, b, c$ adalah bilangan real di interval $(0, 1)$, buktikan bahwa \[\sqrt{abc} + \sqrt{(1 - a)(1 - b)(1 - c)} < 1\]
    % inequalities manfrino ex 1.7.7
\end{problem}

\begin{problem}
    (Latvia MO 2002) Misal $a, b, c, d$ adalah bilangan real positif yang memenuhi \[\dfrac{1}{1 + a^4} + \dfrac{1}{1 + b^4} + \dfrac{1}{1 + c^4} + \dfrac{1}{1 + d^4} = 1\] Buktikan bahwa $abcd \geq 3$
    % Topics in Equalities - Hojoo Lee Chapter 2 2.1 7
\end{problem}

\begin{problem}
    Diberikan bilangan real positif $a, b, c$ yang memenuhi $(a + b)(b + c)(c + a) = 1$. Buktikan ketaksamaan \[ab + bc + ca \leq \dfrac{3}{4}\]
    % cvetkovski 217
\end{problem}
\end{document}

    